\documentclass[11pt]{amsart}
\usepackage[T1]{fontenc}
\usepackage{lmodern}
\usepackage{amsmath,amssymb}
\usepackage{microtype}
\usepackage[hidelinks,pdftitle={Draft: Even subdivisions are contractible}]{hyperref}
\setlength{\emergencystretch}{1em}

\newtheorem{theorem}{Theorem}[section]
\newtheorem{lemma}[theorem]{Lemma}
\newtheorem{corollary}[theorem]{Corollary}

\title[Even subdivisions are contractible]{Even subdivisions are contractible}
\date{Draft, 5 September 2026}

\begin{document}
\begin{abstract}
We prove that a bipartite graph is contractible in the sense of graph
schemes whenever every vertex in one part has degree at most two.
Consequently, replacing every edge of a loopless multigraph by a path of
positive even length produces a contractible graph. The proof projects a
coloured scheme onto a family of graphic matroids. Each edge label occurs
in at most two projections, and a component count verifies the matroid
union inequalities needed for simultaneous disjoint spanning trees.
If the degrees in that part are at most three, we instead obtain a minor
retaining all prescribed roots in the opposite part. Its proof combines
the same packing argument with a strictly decreasing root-relocation
induction. In particular, every $K_{3,n}$ is weakly contractible.
\end{abstract}
\maketitle

\noindent\textit{Status.} Draft based on internally audited proofs;
not externally peer reviewed.

\section{Statement}

All graphs are finite. Target and scheme graphs are simple; auxiliary
multigraphs may have parallel edges but no loops. Following
K\"undgen, Pelsmajer and Ramamurthi~\cite[Definitions 1.1 and 2.1]{KPR},
an \emph{$H$-scheme} in a graph $G$ with $V(H)\subseteq V(G)$ is a family
$(P_{uv}:uv\in E(H))$ of paths satisfying two conditions:
\begin{enumerate}
\item $P_{uv}$ joins $u$ to $v$ and contains no other vertex of $H$;
\item whenever a nonempty collection of scheme paths has a common vertex,
their corresponding edges in $H$ have a common endpoint.
\end{enumerate}
A \emph{rooted $H$ minor} in $G$ is a family $(C_v:v\in V(H))$ of
pairwise disjoint connected vertex sets, with $v\in C_v$ for every $v$,
such that $C_u$ and $C_v$ are adjacent whenever $uv\in E(H)$.
The graph $H$ is \emph{contractible} if every graph containing an
$H$-scheme contains a rooted $H$ minor.

\begin{theorem}\label{thm:main}
Let $H$ be bipartite with a specified bipartition $(A,B)$.
If $d_H(b)\leq2$ for every $b\in B$, then $H$ is contractible.
\end{theorem}

\begin{corollary}\label{cor:even}
Let $F$ be a loopless multigraph. Replace each edge of $F$ by a path of
positive even length, using mutually disjoint new internal vertices.
The resulting simple graph is contractible.
\end{corollary}

\begin{proof}
Place every original vertex, and every even-position internal vertex of
each replacement path, in $A$. Place the odd-position internal vertices
in $B$. Every vertex of $B$ has degree two, so Theorem~\ref{thm:main}
applies. Isolated vertices of $F$ may be retained.
\end{proof}

Theorem~\ref{thm:main} includes $K_{2,n}$ for every $n$, with the part
of order $n$ as $B$. Corollary~\ref{cor:even} includes the graph obtained
by subdividing every edge of any simple graph exactly once. Throughout,
``even'' refers to the length of each replacement path.

\section{Packing spanning trees with shared labels}

\begin{lemma}\label{lem:packing}
Let $J$ be a finite index set. For each $j\in J$, let $M_j$ be a connected
multigraph with a distinguished vertex $r_j$ and an edge set labelled
bijectively by a finite set $E_j$. Suppose that $M_j$ is the union of a
family $\mathcal Q_j$ of paths whose edge sets partition $E_j$, all these
paths contain $r_j$, and every other vertex belongs to at least two of
them. Suppose also that each element of $E=\bigcup_j E_j$ belongs to at
most two sets $E_j$. Then there are pairwise disjoint sets
$T_j\subseteq E_j$ such that $T_j$ is a spanning tree of $M_j$.
Trivial paths are allowed; an empty family is allowed when $V(M_j)=\{r_j\}$.
\end{lemma}

\begin{proof}
Fix $X\subseteq E$. Let $c_j(X)$ count the components, including isolated
vertices, of the spanning subgraph of $M_j$ with edge labels $E_j\cap X$.
Put $m_j=|\mathcal Q_j|$. The component containing $r_j$ meets all $m_j$
paths, and every other component meets at least two. Conversely, a path
$Q$ meets at most $1+|E(Q)\setminus X|$ components: deleting those edges
leaves that many path pieces, each contained in one component. Thus
\[
 m_j+2(c_j(X)-1)
 \leq \sum_C |\{Q\in\mathcal Q_j:V(Q)\cap V(C)\ne\varnothing\}|
 \leq m_j+|E_j\setminus X|.
\]
The last bound uses the partition of $E_j$ into the path
edge sets. It follows that
\begin{equation}\label{eq:components}
 c_j(X)-1\leq\tfrac12|E_j\setminus X|.
\end{equation}

Extend the graphic matroid of $M_j$ to the common ground set $E$ by
declaring all labels outside $E_j$ to be matroid loops. Its rank function
$\rho_j$ satisfies
\[
 \rho_j(E)=|V(M_j)|-1,
 \qquad \rho_j(X)=|V(M_j)|-c_j(X).
\]
Summing~\eqref{eq:components} and using the label multiplicity bound gives
\begin{equation}\label{eq:unionbound}
 \sum_j\bigl(\rho_j(E)-\rho_j(X)\bigr)
 \leq \tfrac12\sum_j|E_j\setminus X|
 \leq |E\setminus X|.
\end{equation}
The finite matroid union rank formula~\cite{Edmonds} states that
\[
 \rho_{\bigvee_j\mathcal M_j}(E)
 =\min_{X\subseteq E}\left(|E\setminus X|+\sum_j\rho_j(X)\right).
\]
Equation~\eqref{eq:unionbound}, together with $X=E$, makes this rank equal
to $\sum_j\rho_j(E)$. Write a maximum independent set of the union as
$\bigcup_j I_j$, where each $I_j$ is independent in $\mathcal M_j$.
Equality throughout
\[
 |\bigcup_j I_j|\leq\sum_j|I_j|\leq\sum_j\rho_j(E)
\]
forces the $I_j$ to be pairwise disjoint bases. A base contains no matroid
loops. Therefore $T_j=I_j\subseteq E_j$ is the required spanning tree.
\end{proof}

\section{Proof of Theorem~\ref{thm:main}}

By the coloured-scheme reduction~\cite[Lemma 3.3]{KPR}, an $H$-scheme
has a rooted minor whose underlying graph $G$ supports a \emph{coloured
$H$-scheme}. This consists of a proper colouring
$f:V(G)\longrightarrow V(H)$ fixing each root, such that $P_{uv}$ uses
only colours $u,v$ and every nonroot has degree at least four. We remove
isolated roots before this reduction and restore them as singleton branch
sets at the end. The graph with no remaining roots is immediate.

The paths alternate colours, every edge belongs to exactly one scheme
path, and each nonroot lies on at least two scheme paths
\cite[Remark 3.2(1), (2), (6), (7)]{KPR}. Set
\[
 A_a=f^{-1}(a)\quad(a\in A),
 \qquad L_b=f^{-1}(b)\quad(b\in B).
\]
If $d_H(b)\leq1$, then $L_b=\{b\}$, since no nonroot of colour $b$ can
lie on two distinct paths incident with $b$. If $d_H(b)=2$, every nonroot
of colour $b$ lies internally on both paths incident with $b$. Put
\[
 E=\bigcup_{b\in B}(L_b\setminus\{b\}),
 \qquad E_a=\bigcup_{b\in N_H(a)}(L_b\setminus\{b\}).
\]
These labels are actual vertices of $G$. Each label belongs to exactly
two sets $E_a$, one for each neighbour of its degree-two root colour.

Construct $M_a$ on vertex set $A_a$ as follows. For each $b\in N_H(a)$
and $x\in L_b\setminus\{b\}$, put an auxiliary edge labelled $x$ between
the two neighbours of $x$ on $P_{ab}$. Both neighbours have colour $a$,
and they are distinct because $P_{ab}$ is a path. Different labels may
give parallel edges. Suppress all such vertices $x$ on $P_{ab}$ and
remove its terminal vertex $b$. This leaves a path $Q_b^a$ in $M_a$ from
$a$ to the $A_a$-neighbour of $b$; it is trivial if $P_{ab}=ab$.
Its edge-label set is exactly $L_b\setminus\{b\}$.

As $b$ ranges over $N_H(a)$, the paths $Q_b^a$ partition $E_a$, cover
$A_a$, and all contain $a$. In particular, $M_a$ is connected. Every
vertex of $A_a\setminus\{a\}$ lies internally on at least two distinct
scheme paths incident with $a$, so it remains on at least two projected
paths. Hence the family $(M_a:a\in A)$ satisfies Lemma~\ref{lem:packing}.
Choose its pairwise disjoint spanning-tree label sets $T_a$.

In the reduced graph $G$, define
\[
 C_a=A_a\cup T_a\quad(a\in A),
 \qquad C_b=\{b\}\quad(b\in B).
\]
These sets are pairwise disjoint: the $A_a$ are different colour classes,
the $T_a$ are disjoint sets of nonroot $B$-colour vertices, and no $T_a$
contains a root. Each tree edge labelled $x$ lifts to the two-edge path
through $x$ in $G$, so $G[C_a]$ is connected. Moreover $a\in C_a$.
For each $ab\in E(H)$, the last edge of $P_{ab}$ joins the root $b$ to
a vertex of $A_a\subseteq C_a$. Thus these branch sets form a rooted
$H$-minor model.

Finally compose this model with the rooted minor reduction. Replace every
vertex in each branch set by its connected preimage in the original graph.
Disjointness, connectivity, root containment and required adjacencies are
preserved. Restoring the isolated roots completes the proof.
The singleton assertion for the $B$ branch sets applies only in the
reduced coloured scheme. \hfill$\square$

\section{Degree three with one part of the roots retained}

An $H$-minor model is \emph{rooted on $A\subseteq V(H)$} if $a\in C_a$
for every $a\in A$; roots outside $A$ need not be retained. A target is
\emph{weakly contractible} if every scheme of it contains an unrooted
minor of it.

\begin{lemma}\label{lem:memberships}
Let $H$ be bipartite with bipartition $(A,B)$. If every nonroot of a
$B$ colour in a coloured $H$-scheme belongs to exactly two scheme paths,
then the underlying graph contains a fully rooted $H$ minor.
\end{lemma}

\begin{proof}
Repeat the projections in Section~3, but include a label $x$ in $M_a$
only when $x$ actually occurs on $P_{ab}$. Each label then belongs to
exactly two projections, even when its target colour has degree greater
than two. For fixed $a$, the projected paths partition the labels, cover
$f^{-1}(a)$, and all contain $a$. Every other vertex belongs to at least
two projected paths. Lemma~\ref{lem:packing} therefore supplies disjoint
spanning-tree label sets. The same branch sets $C_a=f^{-1}(a)\cup T_a$,
$C_b=\{b\}$ give the rooted model: the last edge of each $P_{ab}$ still
witnesses its required contact. Isolated roots are singletons.
\end{proof}

\begin{theorem}\label{thm:degree-three}
Let $H$ be bipartite with bipartition $(A,B)$ and $d_H(b)\leq3$ for every
$b\in B$. Every $H$-scheme contains an $H$ minor rooted on $A$.
Consequently $H$ is weakly contractible.
\end{theorem}

\begin{proof}
Delete isolated target roots from the host and restore them as their
original singleton branch sets at the end. They lie on no scheme path;
an edgeless target is immediate. For the remaining fixed target, use
strong induction on host order. During this argument regard $H$ as an
abstract target with an injective current root map $v\mapsto r_v$.

The reduction in~\cite[Lemma 3.3]{KPR} gives a root-preserving minor $K$
supporting a coloured scheme, with $|V(K)|\leq|V(G)|$. Suppose a nonroot
$x$ of colour $b\in B$ lies on every path incident with $b$. Move the
root of $b$ to $x$ and replace each $P_{ba}$ by its suffix from $x$ to
$r_a$. No other path contains $x$, since its colour is $b$. The new paths
avoid the old root $r_b$, contain no other current root internally, and
inherit the scheme intersection condition from the old paths. They form
an $H$-scheme in $K-r_b$, whose order is strictly smaller than $|V(G)|$.
The induction hypothesis gives a model rooted on $A$, since none of those
roots moved. Its lift to $G$ retains every original $A$ root.

Otherwise, every nonroot of a $B$ colour has exactly two memberships.
Indeed each has at least two; a degree-one colour admits none, a
degree-two colour would have full support, and three memberships at a
degree-three colour would also give the preceding case. Apply
Lemma~\ref{lem:memberships} and lift the resulting model to $G$.
Every recursive step strictly lowers host order, and every lift preserves
the required $A$ roots.
\end{proof}

\section{Scope}

The contractible targets supplied by Corollary~\ref{cor:even} have
unbounded treewidth, since they include a once-subdivision of $K_t$ for
every $t$. This does not assert that contractibility is minor-closed or
that $K_t$ itself is contractible. Theorem~\ref{thm:main} also covers
theta graphs whose three paths all have even length; it does not cover
the case of three odd-length paths. It therefore addresses only part of
the bipartite-theta question in~\cite[Section 8, Question 3]{KPR}.
Theorem~\ref{thm:degree-three} gives a $K_{3,n}$ minor retaining its three
prescribed roots in the part of order three. For $K_{3,3}$ it does not
guarantee retention of all six roots; its unrooted special case is already
\cite[Theorem 5.3]{KPR}. Neither theorem establishes full bipartite
contractibility or a Hadwiger conjecture.

Biswal, Lee and Rao~\cite[Lemma 3.2]{BLR} state a broader unrooted claim
for integral unit flows with bipartite demand of minimum degree two.
We use their intended preprint-v2 intersection convention: paths with
four distinct demand endpoints are vertex-disjoint when the independent
intersection count is zero. The published definition appears to reverse
``intersect'' and ``are vertex disjoint''; the following issue concerns
the intended convention, independently of that apparent typo.

Their supplied prefix construction fails even for the coloured
$K_{2,2}$-scheme on eight distinct vertices with paths
\[
 P_{ij}=a_i\,b'_j\,a'_i\,b_j\qquad(i,j\in\{0,1\}).
\]
Each left prefix stops strictly before its first vertex on a path
incident with another left root, giving $C_{a_i}=\{a_i\}$ and
$C_{b_j}=\{b_j,b'_j,a'_0,a'_1\}$. The right sets overlap. This refutes
the construction, including the containment assertion of their
Lemma~3.6, not the intended existence statement of Lemma~3.2. The graph
does have the rooted model $D_{a_i}=\{a_i\}$,
$D_{b_j}=\{b_j,b'_j,a'_j\}$. Our proofs do not use that broader claim;
publication priority remains to be assessed in light of it.

\begin{thebibliography}{9}
\bibitem{BLR}
P.~Biswal, J.~R.~Lee and S.~Rao,
\emph{Eigenvalue bounds, spectral partitioning, and metrical deformations
via flows}, J.~ACM \textbf{57}(3) (2010), Article 13.
\href{https://doi.org/10.1145/1706591.1706593}{doi:10.1145/1706591.1706593}.
Preprint version used for the intersection convention:
\href{https://arxiv.org/abs/0808.0148v2}{arXiv:0808.0148v2}.

\bibitem{Edmonds}
J.~Edmonds, \emph{Submodular functions, matroids, and certain polyhedra},
in Combinatorial Structures and Their Applications, Gordon and Breach,
1970, 69--87.

\bibitem{KPR}
A.~K\"undgen, M.~J.~Pelsmajer and R.~Ramamurthi,
\emph{Finding minors in graphs with a given path structure},
J.~Graph Theory \textbf{79} (2015), 30--47.
\href{https://doi.org/10.1002/jgt.21812}{doi:10.1002/jgt.21812}.
Primary preprint: \href{https://arxiv.org/abs/1207.6141}{arXiv:1207.6141}.
\end{thebibliography}
\end{document}
